\documentclass[11pt,a4paper]{article}

% ------------------------------------------------------------------ Paquets
\usepackage[utf8]{inputenc}
\usepackage[T1]{fontenc}
\usepackage[french]{babel}
\usepackage{lmodern}
\usepackage[margin=2.3cm]{geometry}
\usepackage{booktabs}
\usepackage{array}
\usepackage{float}
\usepackage{caption}
\usepackage{xcolor}
\usepackage{listings}
\usepackage{pgfplots}
\pgfplotsset{compat=1.17}
\usetikzlibrary{babel}          % evite les conflits entre TikZ et babel francais
\usepackage[hidelinks]{hyperref}

% ------------------------------------------------------------------ Style du code
\definecolor{fondcode}{RGB}{247,247,247}
\definecolor{motcle}{RGB}{0,70,160}
\definecolor{commentaire}{RGB}{0,120,60}
\lstset{
  language=C++,
  basicstyle=\ttfamily\small,
  keywordstyle=\color{motcle}\bfseries,
  commentstyle=\color{commentaire}\itshape,
  backgroundcolor=\color{fondcode},
  frame=single, rulecolor=\color{black!20},
  numbers=left, numberstyle=\tiny\color{black!50},
  breaklines=true, columns=fullflexible,
  showstringspaces=false, tabsize=4
}

% ------------------------------------------------------------------ Page de titre
\title{\textbf{Benchmark du parcours en largeur d'un arbre n-aire}\\[4pt]
       \large Influence du conteneur sous-jacent de la file :\\
       \texttt{std::deque}, \texttt{std::list} et file personnalisée sur \texttt{std::vector}}
\author{Prénom NOM \\ ILISI 2 \\ \small Encadrant : \textit{à compléter}}
\date{\today}

\begin{document}
\maketitle

\begin{abstract}
Ce rapport mesure le temps d'exécution d'un parcours en largeur (BFS) sur un arbre n-aire
selon la file utilisée : \texttt{std::queue} reposant sur \texttt{std::deque} (choix par défaut),
\texttt{std::queue} reposant sur \texttt{std::list}, et une file écrite à la main sur un
\texttt{std::vector}. Les mesures sont faites sans aucun affichage, sur un arbre de 12~nœuds
et sur un arbre de 1\,000\,000 de nœuds, avec quatre essais indépendants. La file sur
\texttt{std::list} est la plus lente (environ $5\times$ sur le petit arbre et $8\times$ sur le
grand). \texttt{std::deque} et la file personnalisée obtiennent des temps proches sur le grand arbre.
\end{abstract}

\noindent\textbf{Code source :} \url{https://github.com/rafaabenn/parcours-largeur-arbre}

\tableofcontents
\newpage

% =====================================================================
\section{Objectif}
% =====================================================================

Le parcours en largeur repose sur une file FIFO. En C++, \texttt{std::queue} n'est pas un
conteneur à part entière mais un \emph{adaptateur} : elle enveloppe un autre conteneur
(par défaut \texttt{std::deque}) et n'en expose que \texttt{push}, \texttt{pop}, \texttt{front}
et \texttt{empty}. L'algorithme ne change pas avec le conteneur, mais les performances peuvent
changer. L'objectif de ce benchmark est de répondre à trois questions :

\begin{enumerate}
  \item Quel conteneur sous-jacent rend le parcours le plus rapide ?
  \item L'écart entre les conteneurs dépend-il de la taille de l'arbre ?
  \item Une file simple écrite sur un \texttt{std::vector} est-elle compétitive face à la STL ?
\end{enumerate}

% =====================================================================
\section{Structures comparées}
% =====================================================================

\subsection{L'arbre n-aire}

Chaque nœud stocke ses enfants dans un \texttt{std::vector<Noeud*>}, ce qui permet un nombre
quelconque d'enfants. L'arbre est défini une seule fois dans \texttt{Arbre.h} et il est partagé
par tous les programmes.

\begin{lstlisting}
struct Noeud {
    int valeur;
    std::vector<Noeud*> enfants;   // nombre quelconque d'enfants
    bool visite;
    Noeud(int v) : valeur(v), visite(false) {}
};
\end{lstlisting}

Deux arbres sont utilisés :
\begin{itemize}
  \item \textbf{Petit arbre} (12 nœuds, 4 niveaux, largeur maximale 6) : l'arbre de test du projet,
        dont les nœuds ont 0, 1, 2, 3 ou 4 enfants ;
  \item \textbf{Grand arbre} (1\,000\,000 de nœuds) : arbre complet où chaque nœud a 3 enfants
        (environ 13 niveaux), construit par \texttt{construireGrandArbre(1000000, 3)}.
\end{itemize}

\subsection{Les trois files}

\begin{table}[H]
\centering
\begin{tabular}{>{\ttfamily}l p{9cm}}
\toprule
\normalfont\textbf{File} & \textbf{Fonctionnement} \\
\midrule
queue<Noeud*, deque<Noeud*>> & Choix par défaut de la STL. Éléments rangés par \emph{blocs}
                                contigus ; une allocation seulement quand un bloc est plein. \\
queue<Noeud*, list<Noeud*>>  & Liste doublement chaînée : \emph{une allocation par} \texttt{push},
                                une libération par \texttt{pop}, maillons dispersés en mémoire. \\
File (personnalisée)                & \texttt{std::vector} caché + indice \texttt{debut} :
                                \texttt{pop()} avance l'indice en $O(1)$ sans rien effacer. \\
\bottomrule
\end{tabular}
\caption{Les trois implémentations de file comparées}
\end{table}

\begin{lstlisting}
class File {
private:
    std::vector<Noeud*> elements;
    size_t debut = 0;
public:
    void push(Noeud* n)  { elements.push_back(n); }
    Noeud* front() const {
        if (empty()) throw std::out_of_range("File::front() sur une file vide");
        return elements[debut];
    }
    void pop() {
        if (empty()) throw std::out_of_range("File::pop() sur une file vide");
        debut++;
    }
    bool empty() const   { return debut == elements.size(); }
};
\end{lstlisting}

% =====================================================================
\section{Protocole de mesure}
% =====================================================================

\subsection{Environnement}

\begin{table}[H]
\centering
\begin{tabular}{ll}
\toprule
\textbf{Élément} & \textbf{Valeur} \\
\midrule
Machine       & HP EliteBook 840 G5 \\
Processeur    & Intel Core i7-8550U @ 1,80 GHz \\
Système       & Windows 11 Pro \\
Compilateur   & g++ 15.2.0 (MinGW-W64) \\
Options       & \texttt{-std=c++17 -O2 -Wall -Wextra} \\
Horloge       & \texttt{std::chrono::steady\_clock} \\
\bottomrule
\end{tabular}
\caption{Environnement de test}
\end{table}

\subsection{Mesurer l'algorithme, pas l'affichage}

Un appel à \texttt{std::cout} coûte beaucoup plus cher que l'extraction d'un nœud de la file.
Chronométrer un parcours qui affiche reviendrait surtout à mesurer l'affichage. Le protocole
sépare donc strictement le calcul et l'affichage :

\begin{enumerate}
  \item \textbf{Aucun affichage dans le parcours mesuré.} Chaque nœud est « traité » en
        ajoutant sa valeur à une somme, que la fonction renvoie.
  \item \textbf{Une somme de contrôle.} La somme est cumulée puis affichée à la fin. Elle empêche
        le compilateur, en \texttt{-O2}, de supprimer une boucle dont le résultat ne
        servirait à rien. Elle permet aussi de vérifier que chaque parcours a bien visité tous les nœuds.
  \item \textbf{Répétitions et moyenne.} Un parcours du petit arbre dure moins d'une microseconde,
        ce qui est trop court pour une mesure isolée. On le répète 100\,000 fois (20 fois pour le
        grand arbre) et on divise le temps total par le nombre de répétitions.
  \item \textbf{Code identique pour les trois files.} La file est un paramètre \texttt{template} :
        exactement le même parcours est compilé pour chaque conteneur.
  \item \textbf{Plusieurs essais.} Le programme complet est lancé 4 fois. On retient la médiane,
        moins sensible que la moyenne aux perturbations du système.
\end{enumerate}

\begin{lstlisting}
template <typename TypeFile>
long long parcoursLargeur(Noeud* racine)
{
    TypeFile file;
    long long somme = 0;
    file.push(racine);
    while (!file.empty()) {
        Noeud* courant = file.front();
        file.pop();
        somme += courant->valeur;          // traitement, sans affichage
        for (Noeud* e : courant->enfants)
            file.push(e);
    }
    return somme;
}

template <typename TypeFile>
double mesurer(Noeud* racine, int repetitions, long long& controle)
{
    auto debut = chrono::steady_clock::now();
    for (int r = 0; r < repetitions; r++)
        controle += parcoursLargeur<TypeFile>(racine);
    auto fin = chrono::steady_clock::now();
    return chrono::duration<double, nano>(fin - debut).count() / repetitions;
}
\end{lstlisting}

\subsection{Vérification par la somme de contrôle}

\begin{itemize}
  \item Petit arbre : $7+9+67+1+28+8+10+29+2+57+40+6 = 264$ par parcours, donc
        $264 \times 100\,000 \times 3 = 79\,200\,000$.
  \item Grand arbre : valeurs $i \bmod 100$, soit $10\,000 \times 4\,950 = 49\,500\,000$ par
        parcours, donc $49\,500\,000 \times 20 \times 3 = 2\,970\,000\,000$.
\end{itemize}

\noindent Les deux valeurs affichées par le programme sont exactement celles-ci à chaque essai :
tous les parcours ont bien visité tous les nœuds.

% =====================================================================
\section{Résultats}
% =====================================================================

Les valeurs sont des temps moyens pour \textbf{un} parcours.

\begin{table}[H]
\centering
\begin{tabular}{lccccc}
\toprule
\textbf{File} & \textbf{Essai 1} & \textbf{Essai 2} & \textbf{Essai 3} & \textbf{Essai 4} & \textbf{Médiane} \\
\midrule
\texttt{deque}       & 0,288 & 0,401 & 0,333 & 0,213 & \textbf{0,311} \\
\texttt{list}        & 1,654 & 2,079 & 1,676 & 1,219 & \textbf{1,665} \\
\texttt{File} personnalisée &0,668 & 0,603 & 0,678 & 0,537 & \textbf{0,636} \\
\bottomrule
\end{tabular}
\caption{Petit arbre (12 nœuds, 100\,000 répétitions) — temps en microsecondes}
\label{tab:petit}
\end{table}

\begin{table}[H]
\centering
\begin{tabular}{lccccc}
\toprule
\textbf{File} & \textbf{Essai 1} & \textbf{Essai 2} & \textbf{Essai 3} & \textbf{Essai 4} & \textbf{Médiane} \\
\midrule
\texttt{deque}       & 24,7  & 21,7  & 16,2  & 15,9  & \textbf{19,0}  \\
\texttt{list}        & 244,0 & 148,2 & 156,9 & 137,1 & \textbf{152,6} \\
\texttt{File} personnalisée &19,9  & 16,5  & 39,9  & 18,1  & \textbf{19,0}  \\
\bottomrule
\end{tabular}
\caption{Grand arbre (1\,000\,000 de nœuds, 20 répétitions) — temps en millisecondes}
\label{tab:grand}
\end{table}

\begin{figure}[H]
\centering
\begin{minipage}{0.48\textwidth}
\centering
\begin{tikzpicture}
\begin{axis}[
    ybar, bar width=22pt, width=\textwidth, height=6cm,
    symbolic x coords={deque,list,perso},
    xtick=data, ymin=0, ymax=2,
    xticklabels={\texttt{deque}, \texttt{list}, personnalisée},
    ylabel={µs / parcours},
    title={Petit arbre (12 nœuds)},
    nodes near coords,
    every node near coord/.append style={font=\footnotesize, /pgf/number format/use comma},
    yticklabel style={/pgf/number format/use comma},
    enlarge x limits=0.25
]
\addplot[fill=blue!45] coordinates {(deque,0.311) (list,1.665) (perso,0.636)};
\end{axis}
\end{tikzpicture}
\end{minipage}
\hfill
\begin{minipage}{0.48\textwidth}
\centering
\begin{tikzpicture}
\begin{axis}[
    ybar, bar width=22pt, width=\textwidth, height=6cm,
    symbolic x coords={deque,list,perso},
    xtick=data, ymin=0, ymax=180,
    xticklabels={\texttt{deque}, \texttt{list}, personnalisée},
    ylabel={ms / parcours},
    title={Grand arbre (1\,000\,000 nœuds)},
    nodes near coords,
    every node near coord/.append style={font=\footnotesize, /pgf/number format/use comma},
    yticklabel style={/pgf/number format/use comma},
    enlarge x limits=0.25
]
\addplot[fill=orange!60] coordinates {(deque,19.0) (list,152.6) (perso,19.0)};
\end{axis}
\end{tikzpicture}
\end{minipage}
\caption{Temps médian d'un parcours selon la file (plus bas = meilleur)}
\label{fig:barres}
\end{figure}

\begin{table}[H]
\centering
\begin{tabular}{lcc}
\toprule
\textbf{File} & \textbf{Petit arbre} & \textbf{Grand arbre} \\
\midrule
\texttt{deque}       & $1{,}0\times$ & $1{,}0\times$ \\
\texttt{list}        & $5{,}4\times$ & $8{,}0\times$ \\
\texttt{File} personnalisée &$2{,}0\times$ & $1{,}0\times$ \\
\bottomrule
\end{tabular}
\caption{Temps relatif par rapport à \texttt{deque} (sur les médianes)}
\label{tab:ratio}
\end{table}

% =====================================================================
\section{Analyse}
% =====================================================================

\subsection{\texttt{std::list} : le coût des allocations}

\texttt{std::list} est la plus lente dans tous les essais : environ $5\times$ plus lente
que \texttt{deque} sur le petit arbre et $8\times$ sur le grand. Chaque \texttt{push} alloue un
nouveau maillon sur le tas et chaque \texttt{pop} le libère : pour le grand arbre, cela représente
un million d'allocations et un million de libérations \emph{par parcours}. De plus, les maillons
sont dispersés en mémoire, ce qui exploite mal le cache du processeur. L'écart grandit avec la taille de
l'arbre, car ces coûts s'accumulent alors que ceux de \texttt{deque} restent amortis.

\subsection{\texttt{std::deque} : le bon choix par défaut}

\texttt{deque} range ses éléments par blocs contigus et n'alloue qu'un bloc de temps en temps.
Elle libère aussi les blocs vidés au fur et à mesure : sa mémoire reste proportionnelle à la
\emph{largeur} de l'arbre, comme le veut l'algorithme. C'est la plus rapide sur le petit arbre
et elle est à égalité sur le grand.

\subsection{File personnalisée : rapide, mais gourmande en mémoire}

Sur le grand arbre, la file personnalisée obtient la même médiane que \texttt{deque} (19,0~ms). Elle
s'appuie sur un seul bloc contigu (\texttt{std::vector}), ce qui est idéal pour le cache. Sur le
petit arbre, elle est environ deux fois plus lente que \texttt{deque}, car le \texttt{vector}
grandit par réallocations successives (1, 2, 4, 8\dots{} éléments) à chaque parcours, et ce coût
fixe pèse lourd quand il n'y a que 12 nœuds.

Elle a surtout un défaut : \texttt{pop()} n'efface rien. Le \texttt{vector} garde tous les
éléments passés par la file, soit une mémoire en $O(n)$ au lieu de $O(\text{largeur})$. Sur le
grand arbre, cela représente un million de pointeurs (8~Mo) au lieu de la largeur du dernier niveau.

\subsection{Fiabilité des mesures}

Les temps varient d'un essai à l'autre (par exemple 15,9 à 24,7~ms pour \texttt{deque} sur le
grand arbre), car Windows exécute d'autres tâches pendant la mesure. L'essai 3 de la file personnalisée
(39,9~ms) en est un exemple : c'est une valeur isolée. C'est pourquoi on compare les \textbf{médianes} et
non un essai unique. Le \emph{classement}, lui, est stable : \texttt{list} est la dernière dans
les 4 essais et sur les deux tailles d'arbre.

% =====================================================================
\section{Conclusion}
% =====================================================================

\begin{itemize}
  \item Le conteneur sous-jacent ne change pas le résultat du parcours, mais il change nettement
        les performances.
  \item \texttt{std::list} est à éviter pour une file : une allocation par élément la rend
        $5$ à $8$ fois plus lente.
  \item \texttt{std::deque}, choix par défaut de \texttt{std::queue}, offre le meilleur compromis
        entre vitesse et mémoire.
  \item Une file personnalisée sur \texttt{std::vector} est aussi rapide que \texttt{deque} sur un grand arbre,
        mais elle consomme $O(n)$ en mémoire. Elle convient pour un parcours unique ou à des fins
        pédagogiques.
  \item Une mesure de temps n'a de sens que sans affichage, répétée, compilée avec optimisation
        et vérifiée par une valeur de contrôle.
\end{itemize}

\paragraph{Perspectives.}
Remplacer la file personnalisée par un buffer circulaire pour ramener sa mémoire à $O(\text{largeur})$.
Mesurer aussi la mémoire maximale utilisée par chaque file. Comparer avec un parcours en
profondeur utilisant \texttt{std::stack}.

\end{document}
